% Unit stabil: chapter9-unit-133; sumber chapter9.tex, baris 1805--1949.
\section{Teorema Tannaka--Krein untuk Grup Hingga}\label{sec:TK}
Teori klasik Tannaka--Krein \cite{Ta38, Kr49} terutama mengkaji pertanyaan
berikut: untuk grup topologis kompak $G$, bagaimana $G$ dapat
direkonstruksi dari kategori representasinya? Jawabannya berkaitan erat dengan
struktur hasil kali tensor dan dualitas; lihat \cite[\S 4.3]{FL14}. Bagian ini
memberikan uraian ringkas teorema rekonstruksi Tannaka--Krein berdasarkan
hasil-hasil sebelumnya, tetapi hanya untuk grup hingga, bukan grup kompak.

Misalkan $G$ suatu grup dan $\Bbbk$ suatu medan. Selanjutnya, tuliskan
$\otimes := \otimes_{\Bbbk}$ dan unsur identitas $G$ sebagai $1_G$.

Objek utama yang kita perhatikan adalah modul-$G$ yang diperkenalkan dalam Definisi \ref{def:G-mod}, yang dalam banyak konteks juga disebut representasi $G$. Singkatnya, objek-objek ini adalah ruang vektor-$\Bbbk$ $V$ yang dilengkapi aksi kiri linear $G \times V \to V$, dengan aksi ditulis sebagai perkalian kiri $(g, v) \mapsto gv$. Semua modul-$G$ membentuk kategori abelian $\Bbbk$-linear $G\dcate{Mod}$, dan $\Hom$ di dalamnya dinotasikan dengan $\Hom_G$. Funktor pelupa
\[ \omega: G\dcate{Mod} \to \cate{Vect}(\Bbbk) \]
tentu saja memetakan modul-$G$ $V$ ke ruang vektor-$\Bbbk$ yang mendasarinya, yaitu $V$.
\index{representasi}

Tuliskan aljabar grup $G$ sebagai $\Bbbk[G]$; menurut
Contoh \ref{eg:group-Hopf-algebra}, aljabar ini merupakan aljabar Hopf.
Sebagaimana dinyatakan dalam Proposisi \ref{prop:G-mod-kG},
$G\dcate{Mod}$ isomorfik dengan $\Bbbk[G]\dcate{Mod}$. Selanjutnya kita akan
beralih dari satu sudut pandang ke sudut pandang lainnya tanpa penjelasan
tambahan.

\begin{convention}
	Dimensi suatu modul-$G$ berarti dimensinya sebagai ruang vektor-$\Bbbk$. Tuliskan $G\dcate{Mod}_{\mathrm{f}}$ untuk subkategori penuh dalam $G\dcate{Mod}$ yang terdiri atas modul-$G$ berdimensi hingga.
\end{convention}

Medan $\Bbbk$ yang dilengkapi aksi $G$ trivial disebut modul-$G$ trivial
(Contoh \ref{eg:trivial-G-mod}). Untuk modul-$G$ $V$ dan $W$, kita dapat
melengkapi $V \otimes W$ dan $\Hom(V, W) := \Hom_{\Bbbk}(V, W)$ dengan
struktur modul-$G$ kanonik. Kasus $\Hom(V, \Bbbk)$ pada konstruksi terakhir
disebut modul-$G$ kontragredien dari $V$, dan juga dinotasikan dengan $V^*$
untuk membedakannya dari dual sebagai ruang vektor. Operasi-operasi ini telah
dijelaskan secara terperinci setelah Definisi \ref{def:HomG}.

\begin{proposition}\label{prop:RepG-Tannakian}
	Operasi hasil kali tensor di atas $(V, W) \mapsto V \otimes W$ memberikan struktur monoidal simetris pada $G\dcate{Mod}$, dengan modul-$G$ trivial $\Bbbk$ sebagai unit, dan funktor pelupa $G\dcate{Mod} \to \cate{Vect}(\Bbbk)$ menjadi funktor monoidal terhadap struktur ini. Subkategori penuh $G\dcate{Mod}_{\mathrm{f}}$ menjadi kategori Tannakian netral, sedangkan funktor pelupa $\omega: G\dcate{Mod}_{\mathrm{f}} \to \cate{Vect}_{\mathrm{f}}(\Bbbk)$ merupakan funktor seratnya.
\end{proposition}
\begin{proof}
	Struktur monoidal simetris pada $G\dcate{Mod}$ dapat diperiksa langsung
	dengan mengikuti definisi hasil kali tensor. Pembaca yang menyukai gambaran
	besar juga dapat memahaminya sebagai berikut.
	\begin{itemize}
		\item Karena $\Bbbk[G]$ merupakan aljabar Hopf, Proposisi \ref{prop:bialgebra-Mod-monoidal} memberikan struktur monoidal pada $\Bbbk[G]\dcate{Mod}$ sedemikian sehingga funktor pelupanya monoidal. Berdasarkan fakta bahwa komultiplikasi pada $\Bbbk[G]$ adalah $g \mapsto g \otimes g$, dengan menguraikan definisi secara teliti terlihat bahwa struktur ini tepat merupakan operasi hasil kali tensor yang didefinisikan sebelumnya.
		\item Karena $\Bbbk[G]$ kokomutatif, Proposisi \ref{prop:bialgebra-mod-symmetry} (i) juga mengakibatkan bahwa $\Bbbk[G]\dcate{Mod}$ merupakan kategori monoidal simetris.
	\end{itemize}
	
	Jelas bahwa $G\dcate{Mod}_{\mathrm{f}}$ merupakan subkategori monoidal dari $G\dcate{Mod}$, dan sekaligus merupakan kategori abelian. Untuk setiap $n \in \Z_{\geq 0}$, semua struktur modul-$G$ yang dapat diberikan pada $\Bbbk^n$ membentuk suatu himpunan kecil; dari sini terlihat bahwa $G\dcate{Mod}_{\mathrm{f}}$ secara esensial kecil. Jelas pula bahwa
	\[ \End_G\left(\Bbbk:\; \text{modul-$G$ trivial}\right) = \Bbbk. \]
	
	Dengan memeriksa langkah demi langkah definisi dualitas dalam \S\ref{def:dualizable-obj}, terlihat bahwa operasi kontragredien $V \mapsto V^*$ menjadikan $G\dcate{Mod}_{\mathrm{f}}$ kategori kaku; data dualitasnya berasal dari morfisme-morfisme yang telah dikenal
	\[\begin{tikzcd}[row sep=tiny]
		V \otimes V^* \arrow[r, "{\mathrm{ev}}"] & \Bbbk \\
		v \otimes \lambda \arrow[mapsto, r] & \lambda(v),
	\end{tikzcd} \quad \begin{tikzcd}[row sep=tiny]
		\Bbbk \arrow[r, "{\mathrm{coev}}"] & V^* \otimes V \\
		1 \arrow[mapsto, r] & \sum_i \check{v}_i \otimes v_i,
	\end{tikzcd}\]
	di mana $(v_i)_i$ adalah sembarang basis bagi $V$ sebagai ruang vektor, dan
	$(\check{v}_i)_i$ adalah basis dualnya; keduanya merupakan morfisme dalam
	$G\dcate{Mod}$. Jika menggunakan sudut pandang aljabar Hopf $\Bbbk[G]$,
	kekakuan juga dapat diturunkan dari Proposisi \ref{prop:Hopf-module-dual}
	dan fakta bahwa antipoda pada $\Bbbk[G]$ adalah $g \mapsto g^{-1}$.
	
	Konstruksi di atas juga mengakibatkan bahwa funktor pelupa $\omega: G\dcate{Mod}_{\mathrm{f}} \to \cate{Vect}_{\mathrm{f}}(\Bbbk)$ merupakan funktor monoidal yang kompatibel dengan struktur kategori monoidal simetris pada kedua sisi, dan jelas pula bersifat eksak. Jadi, $\omega$ merupakan funktor serat netral dari kategori Tannakian netral $G\dcate{Mod}_{\mathrm{f}}$.
\end{proof}

Selanjutnya, misalkan $G$ grup hingga. Pokok perhatian bagian ini adalah:
\begin{center}
	Bagaimana grup $G$ dapat direkonstruksi dari data $\left(G\dcate{Mod}_{\mathrm{f}}, \omega\right)$?
\end{center}
Hasil-hasil yang diperoleh dalam \S\ref{sec:Tannakian-cat} telah mendekati jawaban atas pertanyaan ini, tetapi perangkatnya masih perlu sedikit disesuaikan. Pertama-tama, kita perlu menafsirkan $G\dcate{Mod}_{\mathrm{f}}$ sebagai kategori komodul.

\begin{definition}
	Tetap misalkan $G$ grup hingga. Definisikan ruang vektor-$\Bbbk$
	\[ C(G) := \left\{\text{pemetaan}\; f: G \to \Bbbk \right\}, \]
	dengan struktur ruang vektornya berasal dari operasi titik demi titik. Selain itu, terdapat isomorfisme kanonik
	\[\begin{tikzcd}[row sep=tiny]
		C(G) \otimes C(G) \arrow[r, "\sim"] & C(G \times G) \\
		f_1 \otimes f_2 \arrow[mapsto, r] & {[(g_1, g_2) \mapsto f_1(g_1) f_2(g_2)]}.
	\end{tikzcd}\]
\end{definition}

Setelah $C(G) \otimes C(G)$ diidentifikasi dengan $C(G \times G)$, $C(G)$
memiliki struktur berikut:
\begin{center}\begin{tabular}{|c|c|c|c|c|} \hline
	Perkalian & Komultiplikasi $\Delta$ & Unit & Kounit $\epsilon$ & Antipoda $S$ \\ \hline
	Perkalian titik demi titik & $f \mapsto \left[ (g_1, g_2) \mapsto f(g_1 g_2) \right]$ & Fungsi konstan $1$ & $f \mapsto f(1_G)$ & $(Sf)(g) = f(g^{-1})$ \\ \hline
\end{tabular}\end{center}
Pemeriksaan langsung menunjukkan bahwa struktur ini merupakan aljabar Hopf
komutatif. Oleh karena itu, Proposisi \ref{prop:bialgebra-mod-symmetry} (ii)
menjamin bahwa $\cated{Comod}C(G)$ menjadi kategori monoidal simetris, sedangkan
Proposisi \ref{prop:Hopf-module-dual} menjadikan
$\catesubd{Comod}{\mathrm{f}}C(G)$ kategori kaku monoidal simetris; bandingkan
dengan bagian ``gambaran besar'' dalam pembuktian sebelumnya.

Untuk setiap $g \in G$, definisikan $\mathbf{1}_g \in C(G)$ sebagai unsur yang
bernilai $1$ pada $g$ dan bernilai $0$ di tempat lain. Unsur-unsur ini
membentuk basis bagi $C(G)$, sedangkan perkalian dan komultiplikasinya
masing-masing ditentukan oleh
$\mathbf{1}_{g_1} \mathbf{1}_{g_2} = \delta_{g_1, g_2} \mathbf{1}_{g_1}$
(simbol $\delta$ Kronecker) dan
$\Delta(\mathbf{1}_g) = \sum_{xy=g} \mathbf{1}_x \otimes \mathbf{1}_y$.

\begin{proposition}\label{prop:Rep-comod-fiber}
	Terdapat isomorfisme kategori
	$G\dcate{Mod} \simeq \cated{Comod}C(G)$: komodul kanan $C(G)$ yang
	bersesuaian dengan modul-$G$ $V$ adalah $V$ yang dilengkapi pemetaan linear
	\[\begin{tikzcd}[row sep=tiny]
		\rho: V \arrow[r] & V \otimes C(G) \arrow[r, "\sim"] & \left\{\text{pemetaan}\; G \to V \right\} \\
		v \arrow[mapsto, r] & \displaystyle\sum_{g \in G} gv \otimes \mathbf{1}_g \arrow[mapsto, r] & {[g \mapsto gv]},
	\end{tikzcd}\]
	sedangkan modul-$G$ yang bersesuaian dengan komodul kanan $C(G)$ $(V, \rho)$ adalah $V$ yang dilengkapi pemetaan
	\begin{align*}
		G \times V & \to V \\
		(g, v) & \mapsto \underbracket{\rho(v)}_{\text{pemetaan}\; G \to V}(g).
	\end{align*}

	Lebih lanjut, struktur monoidal simetris pada $G\dcate{Mod}$ dan $\cated{Comod}C(G)$ juga bersesuaian dengan cara ini. Semua isomorfisme tersebut mempertahankan funktor pelupa menuju $\cate{Vect}(\Bbbk)$, sebagaimana dinyatakan oleh diagram komutatif
	\[\begin{tikzcd}[column sep=small]
		G\dcate{Mod} \arrow[rr, "\sim"] \arrow[rd, "\omega"'] & & \cated{Comod}C(G) \arrow[ld, "{\omega_{C(G)}}"] \\
		& \cate{Vect}(\Bbbk). &
	\end{tikzcd}\]

	Secara khusus, $G\dcate{Mod}_{\mathrm{f}}$ dan $\catesubd{Comod}{\mathrm{f}}C(G)$ juga saling isomorfik, dan isomorfisme tersebut mempertahankan funktor pelupa menuju $\cate{Vect}_{\mathrm{f}}(\Bbbk)$.
\end{proposition}
\begin{proof}
	Periksa langsung definisi-definisinya.
\end{proof}

\begin{lemma}\label{prop:CG-group}
	Terdapat bijeksi
	$G \xrightarrow{1:1} \Hom_{\Bbbk\dcate{CAlg}}(C(G), \Bbbk)$ yang memetakan
	$g$ ke homomorfisme evaluasi $f \mapsto f(g)$. Lengkapi ruas kanan dengan
	operasi biner yang berasal dari
	$\Delta: C(G) \to C(G \times G) \simeq C(G) \otimes C(G)$ dan memetakan
	$(\varphi_1, \varphi_2)$ ke
	$(\varphi_1 \otimes \varphi_2) \circ \Delta$. Dengan operasi ini, bijeksi
	tersebut merupakan isomorfisme grup.
\end{lemma}
\begin{proof}
	Setiap homomorfisme aljabar-$\Bbbk$ $\varphi: C(G) \to \Bbbk$ bersifat surjektif, sehingga $\Ker(\varphi)$ merupakan ideal maksimal. Namun, sebagai aljabar-$\Bbbk$, $C(G)$ adalah hasil kali langsung dari $G$ salinan $\Bbbk$, melalui pemetaan yang membawa $f$ ke $(f(g))_{g \in G}$. Selain itu, mudah ditunjukkan bahwa ideal-ideal maksimal dari $\prod_{g \in G} \Bbbk$ berkorespondensi secara bijektif dengan unsur-unsur $G$ melalui
	\[  g \in G \quad \leftrightarrow \quad \Bbbk \times \cdots \times \underbracket{\{0\}}_{\text{komponen $g$}} \times \cdots \times \Bbbk \]
	Pembaca dapat melihat \cite[Lema 8.9.7]{Li1} atau membuktikannya sendiri. Dari sini, bijeksi yang diminta segera diperoleh.
	
	Untuk membuktikan bahwa korespondensi ini mempertahankan perkalian grup,
	cukup diperiksa bahwa
	\[\begin{tikzcd}
		(g_1, g_2) \arrow[phantom, r, "\in" description] \arrow[mapsto, d] & G \times G \arrow[r, "1:1"] \arrow[d] & \Hom_{\Bbbk\dcate{CAlg}}(C(G \times G), \Bbbk) \arrow[d] & {[F \mapsto F(g_1, g_2)]} \arrow[phantom, l, "\ni" description, sloped] \arrow[mapsto, d] \\
		g_1 g_2 \arrow[phantom, r, "\in" description] & G \arrow[r, "1:1"'] & \Hom_{\Bbbk\dcate{CAlg}}(C(G), \Bbbk) & {[f \mapsto (\Delta f)(g_1, g_2)]} \arrow[phantom, l, "\ni" description, sloped]
	\end{tikzcd}\]
	merupakan diagram komutatif. Hal ini tidak menimbulkan kesulitan yang berarti.
\end{proof}

Sekarang kita dapat menjawab pertanyaan rekonstruksi yang diajukan sebelumnya.
Jawabannya melibatkan grup $\Aut^{\otimes}(\omega)$ yang diperkenalkan dalam
Korolari \ref{prop:Aut-Tannakian}; pokok persoalannya ialah mendeskripsikan isomorfisme
tersebut secara eksplisit.

\begin{theorem}[T.\ Tannaka, M.\ Krein]\label{prop:Tannaka-Krein}
	\index{Teorema Tannaka-Krein@Teorema Tannaka--Krein (Tannaka--Krein Theorem)}
	Misalkan $\Bbbk$ suatu medan dan $G$ suatu grup hingga. Tinjau funktor pelupa $\omega: G\dcate{Mod}_{\mathrm{f}} \to \cate{Vect}_{\mathrm{f}}(\Bbbk)$. Maka terdapat isomorfisme grup
	\[ A: G \rightiso \Aut^\otimes(\omega), \]
	yang secara eksplisit dideskripsikan sebagai berikut: jika $g \in G$ dan $V$ suatu modul-$G$ berdimensi hingga, maka $A(g)_V: V \to V$ adalah pemetaan $\Bbbk$-linear $v \mapsto gv$.
\end{theorem}
\begin{proof}
	Diketahui bahwa $\omega$ merupakan funktor serat netral dari kategori
	Tannakian $G\dcate{Mod}_{\mathrm{f}}$
	(Proposisi \ref{prop:RepG-Tannakian}), sehingga diperoleh aljabar Hopf
	komutatif $\coEnd(\omega)$. Dengan mengambil $B = B' = \Bbbk$ dalam
	Korolari \ref{prop:Aut-Tannakian}, diperoleh isomorfisme grup
	\[ \Aut^{\otimes}(\omega) \simeq \Hom_{\Bbbk\dcate{CAlg}}\left(\coEnd(\omega), \Bbbk \right). \]

	Diagram komutatif dalam Proposisi \ref{prop:Rep-comod-fiber} mengakibatkan
	\[ \coEnd(\omega) \simeq \coEnd\left(\omega_{C(G)}: \catesubd{Comod}{\mathrm{f}}C(G) \to \cate{Vect}_{\mathrm{f}}(\Bbbk) \right), \]
	sedangkan Lema \ref{prop:reconstruction-cat-aux} beserta penjelasan sesudahnya memberikan
	\[ \coEnd\left( \omega_{C(G)} \right) \simeq C(G); \]
	semuanya merupakan isomorfisme kanonik bialjabar. Dengan memasukkan Lema \ref{prop:CG-group}, kita memperoleh isomorfisme grup
	\[\begin{tikzcd}[row sep=tiny]
		G \arrow[r, "\sim"] & \Hom_{\Bbbk\dcate{CAlg}}\left(C(G), \Bbbk \right) \arrow[r, "\sim"] & \Aut^{\otimes}(\omega) \\
		g \arrow[mapsto, r] & B(g) \arrow[mapsto, r] & A(g).
	\end{tikzcd}\]
	Citra $B(g)$ di bawah isomorfisme pertama ialah homomorfisme evaluasi
	$f \mapsto f(g)$. Di sisi lain, berdasarkan sifat universal
	$\coEnd(\omega) = \coHom(\omega, \omega) \simeq C(G)$ dalam
	Definisi \ref{def:cogebra-L-universal} (dengan mengambil $L = \Bbbk$),
	isomorfisme kedua sebenarnya meluas menjadi
	\[ \Hom_{\Bbbk}(C(G), \Bbbk) \rightiso \Hom_{\Bbbk}(\omega, \omega), \]
	dan sifat universal tersebut menunjukkan bahwa korespondensi antara $A$ dan $B$ dicirikan oleh diagram komutatif dalam $\cate{Vect}_{\mathrm{f}}(\Bbbk)$:
	\[\begin{tikzcd}[column sep=large]
		& \omega(V) \arrow[ld, "\rho"'] \arrow[rd, "{A(g)_V}"] & \\
		\omega(V) \otimes C(G) \arrow[r, "{\identity \otimes B(g)}"'] & \omega(V) \otimes \Bbbk \arrow[r, "\sim"'] & \omega(V),
	\end{tikzcd}\]
	di mana $V$ berkisar atas semua modul-$G$ berdimensi hingga dan $\rho$
	didefinisikan seperti dalam Proposisi \ref{prop:Rep-comod-fiber}. Dengan
	memasukkan $B(g)(f) = f(g)$, segera terlihat bahwa $A(g)_V(v) = gv$. Ini
	membuktikan pernyataan tersebut.
\end{proof}

Di bawah syarat yang sesuai pada $\Bbbk$, kita dapat mencirikan lebih lanjut
aljabar Hopf komutatif mana saja yang isomorfik dengan $C(G)$ untuk suatu grup
hingga $G$. Pertanyaan ini termasuk dalam ranah geometri aljabar elementer atau
teori gelanggang komutatif dan tidak akan dibahas lebih jauh di sini.

\index{Teorema Doplicher-Roberts@Teorema Doplicher--Roberts (Doplicher--Roberts Theorem)}
Terakhir, terdapat metode lain yang tidak bergantung pada funktor serat dan merekonstruksi grup kompak semata-mata dari kategori modul-$G$, yang disebut teorema rekonstruksi Doplicher--Roberts dan berakar dalam teori medan kuantum. Pembaca yang berminat dapat melihat survei \cite{Mu07}.
